% scaling.tex -- running time as a function of input size.
% Rules shown (Sanders, "Presenting Data from Experiments in Algorithmics"):
%  - log x axis with powers of two, measured at powers of two;
%  - y normalized (time per element) so the interesting part is visible,
%    instead of a log-log plot that is close to a diagonal;
%  - y starts at 0, so a small difference does not look like a large one;
%  - points mark measurements, lines only connect points of one curve;
%  - legend in the empty upper left, in the order of the curves at the
%    right end;
%  - the plot shows where a baseline wins (BaseB for large n). Keep that.
% Build from the figures directory: make build/scaling.pdf
\documentclass[border=2pt]{standalone}
\input{../latex/figure-preamble}
\input{../palette/colors}
\begin{document}
\begin{tikzpicture}
\begin{axis}[
  xmode=log,
  xlabel={input size $n$},
  ylabel={time per element [ns]},
  ymin=0, ymax=21,
  xmin=2^9.5, xmax=2^24.5,
  xtick={1024,16384,262144,4194304},
  legend pos=north west,
  table/col sep=space,
]
  % Order of \addplot = order in the legend = order at the right end.
  \addplot[baseA] table[x=n, y=baseA, comment chars={\#}] {data/scaling.dat};
  \addlegendentry{\algname{BaseA}}
  \addplot[ours]  table[x=n, y=ours,  comment chars={\#}] {data/scaling.dat};
  \addlegendentry{\algname{Ours}}
  \addplot[baseB] table[x=n, y=baseB, comment chars={\#}] {data/scaling.dat};
  \addlegendentry{\algname{BaseB}}
\end{axis}
\end{tikzpicture}
\end{document}
